% Compile with: uplatex file.tex && dvipdfmx file.dvi
%           or: latexmk -uplatex file.tex
\documentclass[dvipdfmx]{article}
\usepackage{amsmath, amssymb}
\usepackage{tikz}

\title{scenario\_hyperfine\_structure}
\date{\today}

\begin{document}
\maketitle

The hyperfine interaction $H_{\rm hf}=A\,\mathbf I\cdot\mathbf S$, taken along two independent routes. $A$ is the hyperfine coupling constant.

Define A as a scalar

\[
  A
\]

[Route A] Couple the nuclear spin $I=3/2$ to the electron spin $S=1/2$. This gives the allowed total angular momenta $F$ and an identity expressing $\mathbf I\cdot\mathbf S$ through $F$. $F$ stays symbolic, so the identity can be carried on by ordinary operations.

couple I=3/2 and S=1/2

\begin{equation}
  \left| \frac{3}{2} \right\rangle \otimes \left| \frac{1}{2} \right\rangle:\quad F = 1,\; 2
  \label{eq:E1}
\end{equation}

\begin{equation}
  \mathbf{I}\cdot\mathbf{S} = \frac{1}{2}\left[ F(F+1) - \frac{15}{4} - \frac{3}{4} \right]
  \label{eq:E2}
\end{equation}

Multiply both sides by the coupling constant $A$ to get the hyperfine eigenvalues — and to check that an identity built without the CAS rides an ordinary both-sides operation.

multiply both sides of E2 by A

\begin{equation}
  A \mathbf{I}\cdot\mathbf{S} = \frac{2 A F^2+2 A F-9 A}{4}
  \label{eq:E3}
\end{equation}

Substitute the upper level $F=2$.

in E3, substitute 2 for F

\begin{equation}
  A \mathbf{I}\cdot\mathbf{S} = \frac{3 A}{4}
  \label{eq:E4}
\end{equation}

Substitute the lower level $F=1$.

in E3, substitute 1 for F

\begin{equation}
  A \mathbf{I}\cdot\mathbf{S} = -\frac{5 A}{4}
  \label{eq:E5}
\end{equation}

Take the interval between the two levels. The Landé interval rule $\Delta E = A\,F_{\rm max}$ says it should be $2A$.

$E4 - E5$

\begin{equation}
  E4 - E5
  \label{eq:E6}
\end{equation}

simplify E6

\begin{equation}
  2 A
  \label{eq:E7}
\end{equation}

[Route B] Now the same $\mathbf I\cdot\mathbf S$ through a spin-algebra trace at infinite temperature. First declare each site's spin — which spins the trace was taken with is part of what the result means, so this is recorded as a timeline step.

site 1 has spin 1/2 and site 2 has spin 3/2

\noindent\textbf{Assumption:}\ $\displaystyle $\par

Register $\mathbf I\cdot\mathbf S$ as the sum of its $x,y,z$ components.

$S^x_1 S^x_2 + S^y_1 S^y_2 + S^z_1 S^z_2$

\begin{equation}
  S^x_1 S^x_2 + S^y_1 S^y_2 + S^z_1 S^z_2
  \label{eq:E8}
\end{equation}

The first moment (expectation value), which should vanish at infinite temperature.

expectation value of E8

\begin{equation}
  \langle E8 \rangle = 0
  \label{eq:E9}
\end{equation}

The second moment. With $\frac13 x_1x_2$ and $x_i=S_i(S_i+1)$, putting in $x_1=3/4$ and $x_2=15/4$ should give $15/16$ — confirmation that the per-site spins are doing the work.

2nd moment of E8

\begin{equation}
  \langle {E8}^{2} \rangle = \frac{15}{16}
  \label{eq:E10}
\end{equation}

\end{document}